\documentclass[11pt,a4paper]{article}

\usepackage[margin=2.5cm]{geometry}
\usepackage[T1]{fontenc}
\usepackage[utf8]{inputenc}
\usepackage{lmodern}
\usepackage{graphicx}
\usepackage{float}
\usepackage{booktabs}
\usepackage{caption}
\usepackage{parskip}
\usepackage[hidelinks]{hyperref}
\usepackage[none]{hyphenat}  % no word is broken across a line
\sloppy

\graphicspath{{out/}}
\captionsetup{font=small,labelfont=bf}

\newcommand{\userA}{\textit{triflicacid}}
\newcommand{\userB}{\textit{megamoleys}}
% A quoted argument and a long option, as typed at the command line.
\newcommand{\q}[1]{\textquotedbl #1\textquotedbl}
\newcommand{\opt}[1]{-{}-#1}
% \chart{file}{caption}{command that produced the file}
\newcommand{\chart}[3]{%
  \begin{figure}[H]
    \centering
    \includegraphics[width=\textwidth,height=0.78\textheight,keepaspectratio]{#1}
    \caption{#2 Generated with: \texttt{#3}}
  \end{figure}}

\title{A Quantitative Analysis of a Two Person Message History}
\author{}
\date{21 September 2025 to 7 October 2026}

\begin{document}
\maketitle

\begin{abstract}
This report analyses a direct message history between two participants, \userA{} and \userB{}.
The history contains 39,072 messages and 212,571
words sent over 382 days. Nineteen charts are presented, each describing one measurable
aspect of the conversation. Overall volume was shared almost equally throughout. The
measures that changed were those of initiative. From May 2026 one participant started
fewer conversations each month and by September seldom sent the first message of the
day or followed up after silence, and the other participant did more of each.
\end{abstract}

\tableofcontents

\section{Chronology}

Table~\ref{tab:chronology} gives the principal dates of the relationship, which are
referred to throughout.

\begin{table}[H]
  \centering
  \small
  \begin{tabular}{ll}
    \toprule
    Date & Event \\
    \midrule
    21 September 2025 & First message in the history \\
    15 November 2025 & The participants met in person for the first time \\
    30 November 2025 & The relationship began \\
    12 and 13 September 2026 & \userB{} sent no messages \\
    14 September 2026 & \userB{} apologised for the silence and asked to meet \\
    19 September 2026 & The relationship ended \\
    7 October 2026 & Last message in the history \\
    \bottomrule
  \end{tabular}
  \caption{Principal dates in the history.}
  \label{tab:chronology}
\end{table}

The participants began as friends, and for the first ten weeks the conversation
concerned shared interests. They met in person on 15 November 2025. The relationship
began on 30 November 2025, during a conversation that lasted through the night. With
1,801 messages, that day is the busiest in the history.

The relationship lasted a little under ten months. On 12 and 13 September 2026
\userB{} sent nothing. On the morning of 13 September \userA{} wrote, in part:

\begin{quote}
``Is everything okay? I'm guessing you just need some space atm. I'm here if you need
me''
\end{quote}

\userB{} replied on the afternoon of 14 September with four messages sent within
forty seconds of one another:

\begin{quote}
``Hi ruben hope you've been alright''

``Very sorry for the radio silence the past few days ive just been busy and thinking
about things a lot''

``Id like to come down to London this weekend for a day to see and talk to you if
that's okay''

``Would it be alright to hold off finalising plans for the 10th October till after we
chat''
\end{quote}

\userA{} agreed at once, replying ``yeah of course, you're welcome whenever'' and
``sure okay'', and then asked ``are you okay though?''. Within ten minutes the
participants had settled on the Saturday of that week.

They met on Saturday 19 September 2026 and the relationship ended that day. The
decision was made by \userB{}. Contact continued afterwards: message volume had been
recovering since late August and rose a little further in the following two weeks.

\section{Data and Method}

The data is an export of a Discord direct message channel covering 21 September 2025
to 7 October 2026. \userA{} sent 19,337 messages containing 110,000 words, and \userB{}
sent 19,735 messages containing 102,571 words. At least one message was sent on 367 of
the 382 days.

A conversation is defined as a run of messages with no silence longer than
two hours, which yields 1,020 conversations. A turn is a run of consecutive messages
by one participant, and a handover is the point at which the other participant
replies.

Most charts plot weekly values. Weekly counts are small in the quieter months, so the
bold line in each chart is a rolling average over four weeks, drawn over a faint line
showing the raw weekly values. Shares and ratios are calculated over the same four
week window. \userA{} is drawn in orange and \userB{} in navy in every chart.

\section{Overview}

\subsection{Highlights}

The highlights card collects headline figures for the whole period, first for the
pair and then for each participant. The longest unbroken run of daily contact lasted
86 days, from 8 October 2025 to 1 January 2026, and the longest silence 3.1 days. The
largest single conversation contained 1,616 messages over 14.3 hours. A typical one
was far smaller, at nine messages and 56 minutes. The busiest hour was 23:00 and the
busiest day of the week was Sunday. \userA{} started 64\% of conversations and sent
the first message on 274 days, against 94 days for \userB{}.

\chart{highlights.png}{Headline figures for the pair and for each participant.}{python highlights.py}

\subsection{Messages per Day as a Calendar}

Each square in the calendar chart is one day, with weeks running left to right and
weekdays top to bottom. Colour gives the number of messages on a logarithmic scale,
and grey marks a day without any. The grey squares are few and scattered across the
year, and the three longest runs of daily contact lasted 86, 63 and 45 days. The
darkest region runs from late November 2025 to February 2026, and from April 2026 the
squares are noticeably paler.

\chart{calendar_heatmap.png}{Messages on each day of the period, on a logarithmic colour scale.}{python calendar\_heatmap.py}

\section{Volume}

\subsection{Messages per Day}

This chart has three panels. The first plots the number of messages sent each day by
each participant on a logarithmic axis. The five busiest days of each participant are
labelled with the count and the date. Every period in which both participants were
quiet, defined as a combined total of no more than one message in a day, is shaded
grey. Periods of two days or more in which only one participant was quiet, sending no
more than one message a day while the other continued to write, are shaded in that
participant's colour. The second panel shows the running total, and the third shows
each participant's share of each day's messages.

\paragraph{Growth and peaks.}
The history begins slowly, with 142 messages in the final ten days of September 2025
and 2,179 in October. Volume rose to 6,137 messages in November, peaked at 8,048 in
December, and eased to 6,741 in January 2026. The median day across the whole period
contained 51 messages, and ten days exceeded 500. The busiest was 30 November 2025,
with 1,801 messages (882 and 919). The second busiest was 18 January 2026, with
1,589. The three days from 24 to 26 December each contained about 800 to 900. On all
of these days the two participants sent nearly the same number of messages.

\paragraph{Decline and partial recovery.}
After January the monthly total fell in every month until May: 4,414 in February,
2,865 in March, 1,262 in April and 915 in May. The running total in the second panel
climbs steeply until February and is almost flat from April. Volume then recovered
somewhat, to 1,316 in June, 1,257 in July, 1,458 in August and 2,190 in September.
One late peak falls on 25 July 2026, when a single long conversation produced 397
messages. It is the only labelled peak in which the participants differed much, with
247 messages from \userB{} and 150 from \userA{}.

\paragraph{Quiet periods shared by both participants.}
Only 23 of the 382 days contained one message or none. They form the seventeen
periods listed in Table~\ref{tab:quiet}, none of which lasted longer than three days.

Nine of the seventeen were days on which the participants were together in person.
The messages on either side of them are travel arrangements and remarks made on
parting. This group includes both of the three day periods, in the first days of
January 2026 and during a camping trip in August, and single days during visits in
February, March, April, May and June. In this history an absence of messages usually
means that the participants were in the same place.

\begin{table}[H]
  \centering
  \small
  \begin{tabular}{lrrl}
    \toprule
    Period & Days & Messages & Context from adjacent messages \\
    \midrule
    1 to 3 January 2026 & 3 & 1 & Together in person \\
    18 to 20 August 2026 & 3 & 2 & Together in person \\
    23 to 24 September 2025 & 2 & 0 & Before daily contact was established \\
    15 to 16 August 2026 & 2 & 0 & Together in person \\
    26 September 2025 & 1 & 0 & Before daily contact was established \\
    7 October 2025 & 1 & 0 & Before daily contact was established \\
    15 February 2026 & 1 & 0 & Together in person \\
    23 March 2026 & 1 & 0 & Together in person \\
    26 May 2026 & 1 & 0 & Together in person \\
    3 July 2026 & 1 & 0 & Not explained by adjacent messages \\
    19 July 2026 & 1 & 0 & Not explained by adjacent messages \\
    4 October 2026 & 1 & 0 & Not explained by adjacent messages \\
    13 February 2026 & 1 & 1 & Together in person \\
    28 March 2026 & 1 & 1 & Not explained by adjacent messages \\
    25 April 2026 & 1 & 1 & Together in person \\
    20 June 2026 & 1 & 1 & Together in person \\
    12 September 2026 & 1 & 1 & One participant silent \\
    \bottomrule
  \end{tabular}
  \caption{Every period in which the combined total did not exceed one message in a
  day, longest first.}
  \label{tab:quiet}
\end{table}

Three further periods fall in the first weeks of the history, before the participants
wrote to each other every day. Four have no explanation in the surrounding messages.
That leaves 12 September 2026, when the only message was sent by \userA{} and \userB{}
sent nothing on that day or the next. No meeting took place, and the relationship
ended a week later.

\paragraph{Quiet periods of one participant.}
\userA{} had six such periods covering 13 days, and \userB{} had seven covering 18
days. Most are the visits already described, seen from one side. During a visit one
participant often sent a single message on a day when the other sent none, which
makes the period slightly longer for one of them. This accounts for the periods in
the middle of December, February and March, and at the end of May. Three periods have
no such explanation. \userA{} was quiet on 12 and 13 July 2026, and \userB{} on 18 and
19 July and again on 12 and 13 September. In both of the periods belonging to
\userB{} no messages at all were sent on either day. No other period of two days
outside a known visit was entirely silent for either participant.

\paragraph{Share of messages.}
The split between the participants stayed close to even throughout. \userA{} sent
between 44\% and 53\% of each month's messages until July 2026, and about 55\% in
August and September. Daily values scatter more widely in the quieter months, since a
day of twenty messages can be divided more unevenly than a day of five hundred.

\chart{messages_per_day.png}{Daily message counts with peaks and quiet periods marked, running totals, and share of messages.}{python messages\_per\_day.py \opt{total} false \opt{axis} log \opt{quiet} 50 \opt{user-quiet} 100 \opt{quiet-threshold} 1}

\subsection{Message Length}

Messages were short throughout, with overall means of 5.7 words for \userA{} and 5.2
for \userB{}. The chart plots the mean for each participant in each week, with the
middle half of individual messages shaded. The weekly means ran five to six words
during the busiest months, fell to about four in the spring of 2026, and lengthened
again to about six over the summer, with the two lines moving together throughout.

\chart{message_length.png}{Mean words per message in each week, with the middle half of messages shaded.}{python message\_length.py}

\subsection{Volume Breakdown}

The number of words written equals the number of conversations, multiplied by the
messages sent in each conversation, multiplied by the words in each message. This
chart plots the total in its first panel and the three factors beneath it.

Almost all of the decline in volume came from the second factor. The number of
conversations stayed between about 15 and 23 per week in every part of the year. The
number of messages each participant sent in a conversation fell from about 50 in
January 2026 to about seven by April. There were 517 conversations in the first half
of the period and 503 in the second, and the median conversation fell from 15
messages lasting 95 minutes to seven messages lasting 38 minutes.

\chart{volume_breakdown.png}{Words per week and the three factors that multiply to give it.}{python volume\_breakdown.py}

\subsection{Conversation Duration}

The previous chart measured a conversation by the number of messages in it. This one
measures it in time, from the first message to the last. The two upper panels show
the median duration of the conversations started by each participant, with the
middle half of conversations shaded, on a logarithmic axis. The lower panel shows the
difference between the two medians in minutes. Only conversations in which both
participants took part are included, which is 686 of the 1,020. Each point is
calculated from all conversations in the preceding four weeks.

The median conversation lasted about five hours in December 2025, about four hours in
January 2026 and a little over three in February. It fell steadily to about one hour
in June and July, the lowest level of the period, and recovered to about two hours in
September. Duration fell by a factor of about five, which is less than the fall in
the number of messages per conversation.

It made almost no difference which participant started a conversation. The overall
medians were 2.0 hours for those started by \userA{} and 1.9 hours for those started
by \userB{}, and the two upper panels have the same shape. The difference in the
lower panel sits slightly above zero for most of the period and has no trend.

\chart{conversation_duration.png}{Median duration of answered conversations by the participant who started them, and the difference between the two.}{python conversation\_duration.py}

\section{Timing}

\subsection{Message Timing}

Both heat maps are scaled by column, each to its own month, which removes the effect
of the fall in volume. The upper one gives the share of each month's messages sent in
each hour of the day, and the lower one the share sent on each day of the week.

Messaging was concentrated in the late evening, and 23:00 was the busiest hour
overall. In November 2025 about a quarter of all messages were sent between 01:00 and
06:00. That share was about 16\% in January and stayed below 5\% from March onward.
Sunday was the busiest day in most months. The two participants are plotted together
because their timing could not be told apart: the differences between them in any
month were no larger than chance would produce at that volume.

\chart{daily_rhythm.png}{Share of each month's messages by hour of day and by day of week.}{python daily\_rhythm.py}

\subsection{Reply Time}

Replies were fast in every part of the year. The overall median was 16 seconds for
\userA{} and 23 seconds for \userB{}, and neither participant slowed appreciably over
the period. Each participant's weekly median is plotted on a logarithmic axis, with
the shaded band covering the middle half of replies.

\chart{reply_times.png}{Median reply time within a conversation, by week.}{python reply\_times.py}

\subsection{Reply Time Distribution and Longer Waits}

Every handover in the history is included here, divided at two hours. The upper panel
shows, for replies given within two hours, the proportion of messages still unanswered
after each length of time. The lower panel counts the waits that
lasted longer than two hours, grouped by duration, and shows how many of those waits
ended with the waiting participant writing again.

Of 15,085 handovers, 14,514 (96.2\%) were answered within two hours, and half of all
replies arrived within 16 to 23 seconds. One reply in ten took longer than about
eleven minutes. Of the remaining 571 waits, \userA{} was the one waiting on 264
occasions and \userB{} on 307. The two curves in the upper panel are almost
identical.

\chart{response_curve.png}{Distribution of reply times within two hours, and the duration of longer waits.}{python response\_curve.py}

\section{Initiative}

\subsection{Conversation Initiation}

The largest change in the data is here. The three panels show the number of
conversations each participant started in each week, each participant's share of those
starts, and each participant's share of the days on which they sent the first message.

Until April 2026 the two participants
started conversations at broadly similar rates, with \userB{} starting 41 in each of
March and April. The figure for \userB{} then fell to 26 in May, 19 in June, 17 in
July and 13 in August. Over the same months the figure for \userA{} rose, and the
share of conversations started by \userA{} went from 56\% in March to 83\% in
September. The first message of the day follows the same course: \userA{} sent it on
about 67\% of days in March and about 90\% of days from June. The decline began
during an examination period for \userB{} in May and did not reverse when that period
ended.

\chart{initiators.png}{Conversations started per week, share of starts, and share of first messages of the day.}{python initiators.py}

\subsection{Conversation Endings}

The final message of a conversation is the one followed by at least two hours of
silence, and therefore by no reply. This chart shows which participant sent it, with
the lower panel giving each participant's share.

Over the whole period \userA{} sent the final message in 614 conversations (60.2\%)
and \userB{} in 406 (39.8\%). The split was close to even through the spring of 2026.
It diverged from late July and reached about 80\% for \userA{} by late September.
Part of this follows directly from the previous chart, since a conversation that
receives no answer is both started and ended by the same participant.

\chart{last_word.png}{Conversations ended per week and each participant's share.}{python last\_word.py}

\subsection{Consecutive Messages}

A value of one in this chart would mean strict alternation. It plots the mean number
of messages each participant sent in a row before the other replied. Both participants
sent about 2.5 messages per turn for most of the year, and their overall means were
almost identical at 2.56 and 2.62. From August 2026 the figure for \userA{} rose to
about 3.5 and that for \userB{} stayed near 2.7, a gap that lasted about six weeks.

\chart{messages_in_a_row.png}{Mean consecutive messages sent before a reply, by week.}{python messages\_in\_a\_row.py}

\subsection{Nudges}

A nudge is defined here as a message sent by a participant more than one hour after
their own previous message, when no reply had arrived in between. The chart shows the
number of nudges sent each week and each participant's share.

\userA{} sent 556 nudges in total and \userB{} sent 171. The monthly figures for
\userB{} fell from 21 in April to 13 in May, seven in June and no more than three in
any month thereafter. Those for \userA{} rose from about 40 in April to 89 in
September.

\chart{nudges.png}{Nudges sent per week and each participant's share.}{python nudges.py}

\subsection{Questions}

Counting question marks alone gives a misleading result for this pair. \userB{}
omitted the question mark from about three questions in four, and a count based on the
mark attributed 71\% of questions to \userA{}. A message is therefore counted here as a
question if it contains a question mark or is worded as a question, for example by
opening with an interrogative word or ending with a tag such as ``is it''. The chart
gives the weekly counts and each participant's share.

With wording included the totals are 1,669 for \userA{} (54.1\%) and 1,415 for
\userB{} (45.9\%). The share of \userA{} stayed between about 50\% and 60\% for most
of the year, and \userB{} asked the larger number in late January and again in July.
The share moved to about 70\% for \userA{} only in September. The number of questions
asked by \userB{} did not fall during the summer.

\chart{questions.png}{Questions asked per week and each participant's share.}{python questions.py}

\section{Content}

\subsection{Links Shared}

\userB{} shared 1,161 links in total and \userA{} shared 781. Until March 2026
\userB{} shared roughly two to three times as many as \userA{}, for example 199
against 57 in February. The balance reversed in June, and by September \userA{} shared
72 links against 21. The chart gives the weekly counts and each participant's share.

\chart{links.png}{Links shared per week and each participant's share.}{python links.py}

\subsection{Reactions}

The export records how many reactions a message received but not who gave them. In a
conversation between two people a reaction on one participant's message will nearly
always have come from the other. The chart counts the reactions received each week by
each participant's messages.

Messages from \userA{} received 1,144 reactions and messages from \userB{} received
622. Both figures peaked between December 2025 and March 2026. After April the
reactions received by \userA{} settled at about six to fifteen per week and stayed at
that level. Those received by \userB{} fell to about one to three per week. Read the
other way round, as a count of who gave the reactions rather than who received them,
\userB{} went on reacting through the summer and \userA{} largely stopped.

\chart{reactions.png}{Reactions received per week by each participant's messages.}{python reactions.py}

\subsection{Topics}

Eleven topics were specified in advance, each defined by a set of words: work,
university, art, the furry community, games, sleep, food, family, seeing each other,
affection and intimacy. A message is counted under a topic if it contains any of that
topic's words, and may be counted under several. The first row of the chart shows, for
each participant separately, the share of their topic mentions that fell under each
topic. The second row is a heat map of the difference between the two participants,
coloured toward whichever gave a topic the larger share. The third row shows the
proportion of each participant's messages that mentioned any listed topic.

The topics cover only part of the conversation: 15.6\% of the messages of \userA{}
and 17.0\% of those of \userB{} mentioned at least one. The first
row also shows shares, so a topic can widen simply because others have narrowed. In
the quietest months each participant made only about a hundred topic mentions, and
single monthly figures from that period are unreliable.

\paragraph{Before the relationship.}
In October and November 2025 the mix for both participants was dominated by shared
interests. For \userB{} art was the largest single topic in October, at 27\% of topic
mentions, followed by the furry community at 16\%. For \userA{} the furry community
was largest at 23\%, with art, work and games each at about a tenth. Affection and
intimacy together accounted for about a tenth of mentions in October and about a
fifth in November. Across the whole period art was mentioned in 303 messages by
\userB{} and 173 by \userA{}, and the art row of the heat map lies on the side of
\userB{} in these weeks.

\paragraph{The first three months of the relationship.}
From December 2025 affection was the largest topic for both participants, at 28\% of
mentions for \userA{} and 35\% for \userB{} in that month. The shared interests of the
autumn each shrank to between 3\% and 7\%. In January 2026 the largest topic for
\userA{} was intimacy, at 27\%, ahead of affection at 22\%. For \userB{} intimacy
reached 9\% in January, with affection at 33\% and rising to 41\% in February. Over
the whole period intimate vocabulary appeared in 374 messages from \userA{} and 130
from \userB{}. Affectionate vocabulary appeared in 1,042 messages from \userB{} and
615 from \userA{}.

\paragraph{The spring.}
Intimate vocabulary declined first. For \userB{} it fell from 9\% of mentions in
January to 1\% in February and did not return. For \userA{} it fell more gradually,
from 27\% in January to 10\% in February, 7\% in March and 2\% in April. Affectionate
vocabulary took a different course for each participant. For \userB{} it remained at
about a third of all topic mentions from March to May (32\%, 33\% and 36\%), and was
still 28\% in June and 26\% in July. For \userA{} it fell earlier, from 23\% in March
to 12\% in May and 7\% in July. The share of messages from \userB{} that mentioned
any topic was highest in these months, at about 20\% from March to May, largely
because of affectionate terms.

\paragraph{Arrangements to meet.}
The topic of seeing each other, which covers travel and visits, belongs mainly to
\userA{}, with 373 messages against 243. It was the largest topic for \userA{} in May
2026, at 39\% of mentions, and remained at 14\% to 17\% in June and July. For \userB{}
it reached 14\% in May and 17\% in June and was otherwise small.

\paragraph{The summer and the end of the relationship.}
From July 2026 the mix for both participants came to resemble that of the previous
autumn. For \userB{} art returned to about 13\% of mentions from July, the furry
community to 13\% and 14\% in August and September, and games reached 16\% in
September. For \userA{} work was the largest topic in August at 18\%, and the furry
community and art also rose, to 14\% and 13\%. Affection from \userB{} fell to 18\% in
August and to 6\% in September. In that final month affection made up 16\% of the
mentions of \userA{}, against 6\% for \userB{}, and the affection row of the heat map
changes colour.

\chart{topics.png}{Topic share for each participant, the difference between them, and overall coverage.}{python topics.py}

\subsection{Distinctive Words}

This chart lists, for each month, the words that were used at least three times more
often in that month than in the remainder of the history. No topics are specified in
advance, and a word is counted once for each message in which it appears.

October 2025 is dominated by the vocabulary of a factory building game. November
includes words associated with flirting and with meeting in person. January 2026
contains sexually explicit terms alongside programming terms. April is characterised
by a second game, May by travel between two cities, August by exercise and costume
making, and September by a card game. June and July yield few distinctive words.

\chart{distinctive_words.png}{The most distinctive words of each month, coloured by strength.}{python distinctive\_words.py}

\subsection{Word Usage}

Eighteen terms were chosen and their use plotted over time, with one panel for each
and a first panel combining them. Eleven are terms common between partners and four
are names, both of which are discussed below. The remaining
three panels show other expressions characteristic of this pair and are not discussed
here.

\paragraph{Terms common between partners.}
The eleven terms are ``love'', ``dear'', ``miss you'', ``boyfriend'', its
abbreviation ``bf'', ``cute'', ``handsome'', the heart symbol, and all words
beginning with ``kiss'', ``cuddl'' and ``hug''. Use of these rose sharply in December
2025, peaked in January and February 2026, and fell steeply in March. The
participants favoured different terms. \userB{} used ``dear'' 366 times against 39
for \userA{}, ``love'' 399 times against 248, and ``miss you'' 51 times against 12.
\userA{} used words beginning with ``kiss'' 132 times against 62 and the heart symbol
70 times against seven. Most of the terms had fallen to almost nothing by April 2026.
The exception is ``dear'', which \userB{} continued to use about four or five times a
week until September. The phrase ``miss you'' does not appear from \userB{} after May.

\paragraph{Usernames and real names.}
The last four panels show how the participants referred to each other by name. Two
count uses of the usernames, each panel combining the full username with its
shortened forms: ``triflic'' and ``triflicacid'' in one, and ``mega'', ``megamoleys''
and ``moleys'' in the other. The remaining two count uses of the participants' real
first names, ``ruben'' for \userA{} and ``theo'' for \userB{}.

Names of either kind were used rarely, about a hundred times in more than 39,000
messages. The usernames belong almost entirely to the weeks before the relationship
began. \userB{} used the username of \userA{} 24 times, 19 of them in November 2025,
and \userA{} used the username of \userB{} 29 times, 16 of them in the same month.
After November both usernames almost disappear. The real names were used less often
still. \userB{} wrote ``ruben'' 14 times, six of them in December 2025, and \userA{}
wrote ``theo'' six times. \userB{} used ``ruben'' three times in September 2026,
including in the opening message of 14 September quoted in the chronology, having
last used it in March.

Taken together, the eighteen terms were used 1,693 times by \userB{} and 945 times by
\userA{}.

\chart{word_usage.png}{Weekly use of eighteen terms, comprising terms common between partners, three other expressions, and the usernames and real names of the participants.}{python word\_usage.py \opt{rolling} 4 \opt{words} love dear \q{miss you} boyfriend bf cute handsome \q{kiss*} \q{cuddl*} \q{hug*} \q{<3} \q{:3} \q{hai+} meow \q{triflic|triflicacid} \q{mega|megamoleys|moleys} theo ruben}

\section{Qualifications}

Several results do not fit a single account, and some of the measures are weaker than
the charts make them look.

Some measures did not move against \userB{} at all. Reply times were unchanged to the
end. The number of questions asked by \userB{} held up through the summer, and so did
the reactions that \userB{} gave. Conversations started by either participant lasted
equally long.

Intimate vocabulary came mostly from \userA{}, at nearly three times the number of
messages, so its disappearance is as much a change in what \userA{} wrote as in what
\userB{} wrote. The affectionate share of the topic
mentions of \userA{} also fell before that of \userB{}, from 23\% in March to 7\% in
July. Some of that fall is an effect of working in shares, because messages about
arranging visits took up 39\% of the topic mentions of \userA{} in May, though the
data does not show how much of the fall that accounts for.

The timing of the decline in initiative has an obvious alternative explanation. May
2026 was an examination period for \userB{}, and a fall in conversation starts during
examinations would be unremarkable. The data cannot separate that from a change in
the relationship. All that can be said is that the figures did not recover in June,
July or August.

Contact outside the direct message text record is not measured.

Several definitions were chosen for convenience and have not been tested against
alternatives. A conversation ends after two hours of silence, and a different
threshold would change the counts of starts and endings.

\section{Conclusion}

Between late November 2025 and February 2026 the participants exchanged several
thousand messages each month, often through the night. By April the monthly volume
had fallen by more than four fifths. Conversations were held as often as before, and
the fall came from each one being much shorter. Across the whole history the
participants sent almost equal numbers of messages and replied at the same high
speed.

The measures that diverged were those of initiative. \userB{} started progressively
fewer conversations from May 2026. By September more than 80\% of conversation starts
were \userA{}'s, as were almost all nudges. \userA{} also arranged the visits, and was
the participant most often left without a reply. The content of the messages had
changed earlier and in stages. Intimate vocabulary declined from February, first for
\userB{} and then for \userA{}. Affectionate language from \userB{} stayed prominent
through the spring and early summer, declined over July and August, and fell sharply
in September. Questions and reactions from \userB{} continued until the end.

The record therefore shows three changes in sequence rather than one: the vocabulary
of the messages, then their number through the spring, then the initiative behind them
from May. Conduct within a conversation changed least, with replies staying fast to
the end and conversations running the same length whoever had begun them. The
relationship ended on 19 September 2026, and message volume rose slightly in the two
weeks that followed.

\end{document}
